% =============================================================================
\chapter[Partial Differential Equations and FEM]{Partial Differential Equations and the Finite Element Method}
\label{ch:fem_intro}
% =============================================================================

This chapter introduces the continuum problems that motivate the finite element
method, establishes notation used throughout the book, and explains why
Galerkin discretizations provide a systematic route from PDE models to
algebraic systems.

\section{A one-dimensional prototype problem}
\label{sec:fem_intro:prototype}

% A scalar boundary-value problem used to introduce trial spaces, test spaces, and element matrices

\section{Model problems in diffusion, elasticity, and flow}
\label{sec:fem_intro:model}

% Representative scalar, vector, steady, and transient PDEs

\section{Classification of PDEs}
\label{sec:fem_intro:classification}

% Elliptic, parabolic, and hyperbolic behavior; linear vs. nonlinear operators

\section{Boundary and initial data}
\label{sec:fem_intro:data}

% Essential, natural, mixed, and time-dependent data

\section{From strong form to weak form}
\label{sec:fem_intro:strong_to_weak}

% Motivation for test functions, integration by parts, and reduced regularity

\section{From weak form to finite-dimensional approximation}
\label{sec:fem_intro:weak_to_discrete}

% Trial spaces, test spaces, basis functions, and discrete linear systems

\section{Hand assembly of a simple mesh}
\label{sec:fem_intro:hand_assembly}

% A two-element example showing local stiffness matrices, load vectors, and assembly into a global system

\section{Why finite elements?}
\label{sec:fem_intro:motivation}

% Geometry handling, locality, adaptivity, and multiphysics coupling

\section{Roadmap for Part I}
\label{sec:fem_intro:roadmap}

% How the subsequent chapters build the mathematical theory used later
